\chapter{Implementation}
\label{ch:implementation}

This chapter details the implementation of each INR architecture for 3D occupancy prediction on the Nefertiti bust mesh. We discuss the overall code organization, highlight key adaptations required for the 3D setting, and present implementation details for each method. The complete, documented codebase is publicly available at \url{https://github.com/AhmedKElshahed/INR}.

\section{Code Organization}

The implementation follows a modular structure designed for clarity and extensibility. Architecture definitions live in the \texttt{src/models/} package, with each method implemented as a self-contained PyTorch module in its own file (\texttt{siren.py}, \texttt{wire.py}, \texttt{finer.py}, and so on) and exposed through a common \texttt{create\_model} factory. Configuration management is handled by \texttt{src/config.py}, which defines the hyperparameter values used for the 3D experiments alongside the grid-search ranges from which they were selected. Dataset loading and coordinate normalization reside in \texttt{src/data\_3d.py}, and helper functions for visualization and metric computation in \texttt{src/utils.py}.

The main entry point for our experiments is \texttt{train\_3dv2.py}, which orchestrates training, evaluation, and mesh extraction on the Nefertiti bust, and appends one row per model to a results file. A single method can be run in isolation with the \texttt{--models} flag. Dataset generation from raw meshes is handled by \texttt{generate\_data.py}, which performs point sampling and ground truth computation as described in Chapter~\ref{ch:methodology}. Note that \texttt{generate\_data.py} has since been revised, so re-running it produces a different sample distribution than the one used here; the exact dataset behind our results is therefore distributed with the repository as \texttt{nefertiti\_dataset.npz} and should be used to reproduce them.

\section{Adapting 2D Methods for 3D Occupancy}

All eight methods were originally designed and evaluated primarily on 2D tasks such as image reconstruction and novel view synthesis. Adapting them to 3D occupancy requires several modifications that we applied consistently across all architectures.

The most fundamental change is the input dimension, which increases from 2 (for $(x, y)$ coordinates) to 3 (for $(x, y, z)$ coordinates). This affects the first layer's weight matrix dimensions and, for methods using positional encoding, the encoding matrix dimensions. For methods like SIREN with specialized weight initialization, the initialization formulas must be updated to account for the new input dimension.

The output dimension changes from 3 (RGB color values) to 1 (occupancy probability). This simplifies the final layer but also changes the optimization landscape, as the network now solves a binary classification problem rather than continuous regression.

The output activation changes from identity or bounded activations (for RGB in $[0, 1]$) to sigmoid, ensuring outputs represent valid probabilities. The loss function correspondingly changes from Mean Squared Error to Binary Cross-Entropy.

For volumetric evaluation, we generate dense 3D coordinate grids covering the normalized $[-1, 1]^3$ bounding box. The complete grid generation and chunked evaluation logic for memory efficiency is available in the repository.

\section{Architecture Implementations}

The mathematical form of each activation is given in Chapter~\ref{ch:background} and the full hyperparameter values in Appendix~\ref{tab:full_hyperparams}. This section covers only what the 3D implementation required beyond a direct transcription of those definitions.

\subsection{SIREN and FINER}

Both rely on their initialization scheme rather than on any special training machinery. SIREN's first layer uses uniform initialization scaled by the input dimension, while subsequent layers scale by $\omega_0$ to keep activation magnitudes from exploding or vanishing; moving from two input dimensions to three requires updating these formulas accordingly. FINER inherits the same scheme and adds no further implementation concerns, since its variable-periodic activation is a drop-in replacement.

\subsection{WIRE}

WIRE is the only architecture requiring non-trivial changes to the training loop, because its Gabor wavelet activations are complex-valued. PyTorch's autograd propagates complex gradients, but the optimizer expects real parameters, so we take the real part of the gradient before each update. The final layer must also project from complex to real, which we do by taking the magnitude of the complex output before applying the sigmoid.

These two details account for WIRE's cost: storing each weight as a real--imaginary pair gives it 1{,}057{,}282 parameters against SIREN's 264{,}449 at identical depth and width, and it trains 49\% slower than any other method.

\subsection{INCODE}

INCODE's harmonizer is a second network evaluated on every forward pass, predicting the four modulation parameters $a, b, c, d$ from the input coordinates. Despite this, the overhead is modest in practice: INCODE trains 2.6\% slower than Fourier Features on the Nefertiti bust, since the harmonizer is small relative to the main network.

\subsection{Fourier Features}

Fourier Features transform input coordinates through a random projection before feeding them to a standard ReLU network. The projection uses a fixed random matrix $\mathbf{B} \in \mathbb{R}^{3 \times m}$, sampled once at initialization from $\mathcal{N}(0, \sigma^2)$ and held constant thereafter:
\begin{equation}
\gamma(\mathbf{x}) = [\cos(2\pi\mathbf{B}\mathbf{x}), \sin(2\pi\mathbf{B}\mathbf{x})]^\top
\end{equation}

Two distinct parameters control this encoding, and conflating them is an easy mistake to make. The number of random frequencies $m$ determines the width of the encoded input ($2m$ features) and hence the size of the first linear layer; we use $m = 256$, giving a 512-dimensional encoding. The scale $\sigma$ determines how high the sampled frequencies reach; we use $\sigma = 1.0$. Only $\sigma$ controls the frequency content, and Chapter~\ref{ch:results} shows that setting it too high is the single most damaging configuration error we encountered.

\subsection{GAUSS and MFN}

\textbf{GAUSS} uses Gaussian activations $\sigma(x) = \exp(-(sx)^2)$ that provide strong spatial localization. For the Nefertiti bust, we use \texttt{input\_scale = 20.0} and \texttt{$\sigma$ = 1.0}. Notably, GAUSS required a learning rate override of $1 \times 10^{-3}$ (10$\times$ the base rate) to achieve stable convergence, reflecting optimization challenges with this activation function.

\textbf{MFN} (Multiplicative Filter Networks) replaces additive layer composition with element-wise products, creating a fundamentally different computation graph. The multiplicative structure aims to create rich frequency interactions but, as our results show, leads to pathological behavior for occupancy prediction. We use \texttt{input\_scale = 10.0}, \texttt{$\alpha$ = 6.0}, and \texttt{$\beta$ = 1.0}. MFN also required a learning rate override of $3 \times 10^{-4}$ for stable training.

\subsection{FR}

FR (Fourier Reparameterized Training) modifies the weight matrices themselves rather than activations or inputs. Weights are parameterized through learnable coefficients applied to a fixed Fourier basis, encouraging the network to develop frequency-aware representations. We use \texttt{$F$ = 64} frequencies. While conceptually elegant, this approach proved less effective than direct activation modifications for the Nefertiti bust task.

\section{Training Configuration}

The training loop is shared across all methods, with method-specific handling only for gradient processing (WIRE) and loss computation. Key training configuration parameters for the Nefertiti bust experiments:

\begin{itemize}
    \item \textbf{Optimizer}: Adam
    \item \textbf{Epochs}: 500 (all methods)
    \item \textbf{Batch size}: 16,384 points
    \item \textbf{Base learning rate}: $1 \times 10^{-4}$ with cosine annealing to $1 \times 10^{-6}$
    \item \textbf{Learning rate overrides}: GAUSS uses $1 \times 10^{-3}$, MFN uses $3 \times 10^{-4}$
    \item \textbf{Gradient clipping}: Maximum norm 1.0
    \item \textbf{Dataset split}: 450,000 training points, 50,000 validation points
    \item \textbf{Loss function}: Binary Cross-Entropy
    \item \textbf{Output activation}: Sigmoid
\end{itemize}

The dataset is split into training and validation sets to monitor generalization during training. All random seeds are fixed at initialization to ensure reproducibility.

\section{Evaluation Implementation}

Evaluation proceeds in two stages that operate on different data, and the distinction matters when interpreting the results.

\textbf{Evaluation IoU} is computed on the held-out validation split: the 50,000 labelled points withheld from training. Predictions are thresholded at $\tau = 0.5$ and compared against the ray-cast ground truth labels already stored in the dataset. Because the split is produced by a seeded generator over a fixed dataset, every method is evaluated on exactly the same points. We record the final training IoU alongside it, since the difference between the two is what distinguishes a method that has learned the geometry from one that has memorized the training points.

\textbf{Chamfer-L1 and normal consistency} are computed from an extracted surface rather than from point labels. We query the network on a dense $256^3$ grid (16,777,216 points), processed in chunks to stay within GPU memory, and run Marching Cubes at the 0.5 level set to obtain a mesh. Points are then sampled on both this mesh and the ground-truth mesh, and the two metrics are computed from nearest-neighbour correspondences using a k-d tree. These surface samples are drawn at random and are not seeded, so both metrics carry run-to-run variation; differences in the fourth decimal place should not be treated as meaningful.

\section{Reproducibility}

To support reproducibility of our Nefertiti bust experiments:
\begin{itemize}
    \item The dataset is fixed and distributed with the code, so all methods train and are evaluated on identical points
    \item The train/validation split is produced by a seeded generator (\texttt{manual\_seed(42)}), making the 50,000-point held-out set identical across methods and across re-runs
    \item The Nefertiti mesh is normalized to $[-1, 1]^3$ before sampling
    \item Point sampling uses the same strategy (50\% uniform, 50\% near-surface at $\sigma = 0.01$) for all methods
    \item Hyperparameter values are documented in Appendix~\ref{tab:full_hyperparams}
    \item Per-epoch training curves are written to \texttt{logs\_3d/}, and final metrics are appended to \texttt{results\_3d\_comparison.csv}, from which the result tables in this report are generated directly by \texttt{make\_tables.py}
\end{itemize}

One caveat should be stated plainly: only the data split is seeded. Model weight initialization, batch shuffling, and the surface sampling used for Chamfer-L1 and normal consistency all draw from PyTorch's and NumPy's global random state, which the 3D training script does not seed. Repeating a run therefore reproduces the protocol but not the exact numbers. Since each configuration was run once, we cannot separate genuine differences between methods from this run-to-run variation, and we treat small gaps in the results accordingly.

The complete implementation, including all architecture definitions, training scripts, configuration files, and evaluation utilities, is available at \url{https://github.com/AhmedKElshahed/INR}. The repository includes detailed README documentation and instructions for reproducing all experiments reported in this thesis.